% Source for the figure of JacobsonTree (sec. 22.8): a binary tree of 6
% nodes completed with a leaf (grey) in place of each missing child, the
% places of its nodes by levels from 1, and its code B of 2n + 1 bits.
% Build: make figures (in ../build).
\documentclass[border=4pt]{standalone}
\usepackage[T1]{fontenc}
\usepackage{lmodern}
\usepackage{amsmath}
\usepackage{tikz}
\usepackage{lamalunga}
% tikz-ads.tex: the TikZ styles shared by the figures of the memory of the
% manual, \input by each standalone figure: a cell of a vector (cell), the
% cell that an operation touches (touched), a free cell (free), a node of a
% linked structure (node), the labels of indices (index) and of variables
% (var), a pointer (ptr), and \cells and \indices to draw a row of cells.
\usetikzlibrary{arrows.meta,positioning,calc,fit,shapes.multipart}
\definecolor{adsTouched}{RGB}{214,196,236}
\definecolor{adsFree}{RGB}{236,236,236}
\tikzset{
  ads/.style={>={Stealth[length=2mm]},font=\small,
              every node/.append style={inner sep=0pt}},
  cell/.style={draw,minimum width=7mm,minimum height=7mm,
               font=\small\ttfamily,outer sep=0pt},
  touched/.style={cell,fill=adsTouched},
  free/.style={cell,fill=adsFree,text=black!45},
  node/.style={draw,rounded corners=1pt,minimum width=8mm,
               minimum height=6mm,font=\small\ttfamily},
  index/.style={font=\scriptsize\ttfamily,text=black!60},
  var/.style={font=\small\itshape},
  ptr/.style={->,thick}
}
% \cells{name}{x,y}{list}: a row of cells, the first at (x,y), named
% name-0, name-1, ...; each item is text/style, the style one of cell,
% touched and free (e.g., {a/cell,e/touched,/free})
\newcommand{\cells}[3]{%
  \foreach \v/\s [count=\i from 0] in {#3}
    \node[\s] (#1-\i) at ($(#2)+(\i*7mm,0)$) {\v};}
% \indices{name}{n}: the indices 0 .. n-1 under the cells of name
\newcommand{\indices}[2]{%
  \foreach \i in {0,...,\numexpr#2-1\relax}
    \node[index,below=1.2mm of #1-\i] {\i};}

\begin{document}
\begin{tikzpicture}[ads,nd/.style={vertex,minimum size=6.5mm},
    lf/.style={rounded corners=1pt,fill=adsFree,minimum size=3.6mm},
    pl/.style={font=\scriptsize,text=adsInk!60}]
  \node[nd] (A) at (0,0) {A};
  \node[nd] (B) at (-18mm,-11mm) {B};
  \node[nd,vertex-touched,minimum size=6.5mm] (C) at (18mm,-11mm) {C};
  \node[nd] (D) at (-26mm,-22mm) {D};
  \node[lf] (b5) at (-10mm,-22mm) {};
  \node[nd,vertex-touched,minimum size=6.5mm] (E) at (10mm,-22mm) {E};
  \node[nd,vertex-touched,minimum size=6.5mm] (F) at (26mm,-22mm) {F};
  \foreach \n/\x in {8/-31,9/-21,10/5,11/15,12/21,13/31}
    \node[lf] (b\n) at (\x mm,-32mm) {};
  \draw[edge] (A)--(B) (A)--(C) (B)--(D) (B)--(b5) (C)--(E) (C)--(F)
        (D)--(b8) (D)--(b9) (E)--(b10) (E)--(b11) (F)--(b12) (F)--(b13);
  \foreach \n/\p in {A/1,C/3,F/7}
    \node[pl,anchor=south west] at (\n.north east) {\p};
  \foreach \n/\p in {B/2,D/4,E/6}
    \node[pl,anchor=south east] at (\n.north west) {\p};
  \foreach \n/\p in {b5/5,b8/8,b9/9,b10/10,b11/11,b12/12,b13/13}
    \node[pl,anchor=north] at ([yshift=-0.8mm]\n.south) {\p};
  % the code by levels
  \node[var,anchor=east] at (45mm,-11mm) {$B$};
  \foreach \v/\s [count=\i from 0] in {1/cell,1/cell,1/touched,1/cell,0/free,%
      1/touched,1/touched,0/free,0/free,0/free,0/free,0/free,0/free}
    \node[\s,minimum width=6mm] (c\i) at (52mm+\i*6mm,-11mm) {\v};
  \foreach \i in {0,...,12}
    \node[index,below=1.2mm of c\i] {\pgfmathtruncatemacro\p{\i+1}\p};
  \node[var,anchor=west] at (49mm,-24mm)
    {the children of the $x$-th node at $2 x$ and $2 x + 1$:};
  \node[var,anchor=west] at (49mm,-30mm)
    {C, the 3rd node, has E and F at the places 6 and 7};
\end{tikzpicture}
\end{document}
